\section{Guards and Threat Model}
\label{sec:background}

\paragraph{Guards.}
Shieldstral-1.0-3B~\citep{calvi2026shieldstral} uses a fixed system prompt and a user message with three plain-text fields, \texttt{<Instruct>}, \texttt{<Query>} and \texttt{<Document>} (\cref{fig:anatomy}). % src: src/prompt.py (verbatim from the model card)
The developer writes the first two, which together form the deployment's \emph{policy}, and interpolates the content into the third.
Its paper describes no special-token delimiters and reports no prompt-injection evaluation.
We read its score as $P(\text{yes})$, the yes/no logits renormalised over the two answers.
Llama Guard~\citep{inan2023llama,grattafiori2024llama,chi2024llama,meta2025llamaguard4} wraps the conversation in plain-text \texttt{<BEGIN CONVERSATION>} and \texttt{<END CONVERSATION>} lines and has no policy in the prompt for content to override.
gpt-oss-safeguard-20B~\citep{openai2025gptosssafeguard} was post-trained to follow an instruction hierarchy~\citep{wallace2024instruction}, so the policy in its system role sits behind a role boundary the content cannot cross, the design one would pick as a mitigation.

\paragraph{Deployment.}
\label{sec:threat}
A developer places a guard $G$ in front of a downstream system, chooses the policy and the format, and inserts untrusted content $c$ into the \emph{content slot} (Shieldstral's \texttt{<Document>} field).
$G$ \emph{clears} $c$ if its score is below a threshold $\tau$ ($\tau = 0.5$ unless stated).
For \emph{label-only} guards, which return a verdict and no score, $G$ clears $c$ if the verdict is safe.
Cleared content is forwarded unchanged.
Content can arrive as a user message, a retrieved passage, a tool result in an agent loop, or an earlier turn of a conversation.

\paragraph{Attacker.}
The attacker controls only $c$.
They do not control the instruction, the query, the system prompt, decoding or the threshold, and they do not see scores.
For the attack and the benchmark (\cref{sec:attack,sec:benchmark}) the attacker is black-box and may know the deployment's query exactly or not at all.
For the defense (\cref{sec:evaluation}) the attacker knows the defense.
They have the defended weights, the boundary mechanism and the training recipe, may query the defended guard's score, and may optimise content with gradients.

\paragraph{Goal and scope.}
The attacker wants a harmful \emph{request} $h$, a harmful text from a public benchmark that $G$ flags on its own, to be cleared while it stays present and readable in the forwarded content.
Attacks that remove, paraphrase or obfuscate $h$ are out of scope, which separates target rebinding from evasion that hides harmful content from a classifier~\citep{hackett2025bypassing,mangaokar2024prp}.
Attackers who control the guard's prompt or weights, multimodal content and denial of service are also out of scope.

\paragraph{Success criterion.}
Only \emph{screened} requests count, meaning those the deployment flags with no attack.
The screen is $G(h) > 0.9$ in the Shieldstral confirmatory study, a majority \texttt{unsafe} verdict over three samples for label-only guards, and $G(h) \ge 0.5$ in \guardrebind{} and the defense evaluation, so rates from the confirmatory study and the benchmark are not directly comparable. % src: PROTOCOL.md (v1), PROTOCOL_v2.md (Metrics)
The \emph{attack payload} $p$ is the text the attacker adds to $h$, and the attack succeeds if $G$ clears $h \oplus p$, where $\oplus$ is concatenation.
The \emph{bypass rate} (or \emph{bypass}) is the fraction of screened requests on which the attack succeeds, with a Wilson 95\% interval.
We measure success at the guard, not whether a downstream model acts on the request.
